\section{Photometry step by step}

\subsection{Point sources}

\textbf{Step 1 — declare the geometry.} Set \textsf{Source geometry} to
\code{point}. In this mode the program treats your target as an
unresolved star whose light on the detector is spread by the atmosphere
into the seeing profile you selected (Gaussian or Moffat). Everything
that follows --- how much light your aperture captures, how bright the
central pixel gets --- is computed from that profile.

\textbf{Step 2 — choose the aperture radius.} This is the radius, in
arcseconds, of the circular software aperture you will later use in your
photometry program (AstroImageJ, ASTAP, IRIS, \dots); the ETC mirrors it.
The choice is a genuine trade-off. A large aperture captures nearly all
the star's light but also more sky background and more pixels' worth of
read noise; a small one keeps the noise down but loses starlight in the
seeing wings --- and with a Moffat profile the wings carry more light
than intuition suggests. As a rule of thumb, a radius of $1.5\times$ the
seeing FWHM captures about 95\% of a Moffat profile; but you do not need
to optimise by hand, because the ETC applies the exact encircled-energy
correction for whatever radius you type, so every choice gives a correct
(if not optimal) prediction. If you want the optimum, try a few values
and watch the S/N.

\textbf{Step 3 — enter the brightness.} Type the target magnitude, and
verify the two selectors that give that number its meaning: the
\emph{reference filter} (the band your magnitude is expressed in) and the
\emph{magnitude system} (Vega or AB). These need not match the filter you
observe through: it is perfectly normal to say ``the target is $V=12.3$''
while observing in R --- the program derives the colour correction from
the template's spectrum. What matters is that the magnitude, band and
system agree with wherever you got the number from (catalogues usually
state Vega magnitudes for Johnson bands and AB for Sloan bands).

\textbf{Step 4 — run and read.} Press \textsf{Run ETC}. The headline
number is the S/N at your selected time, but look at the whole
S/N-versus-time curve in the plot window: it shows the observation
degrading as the target sets and the airmass grows, and brightening sky
as the Moon rises. The Results window details every contribution
(Sect.~\ref{sec:results}).

One behaviour surprises most users the first time: make the star very
bright and the S/N stops improving, saturating at a few thousand
regardless of magnitude. This is not a bug --- it is atmospheric
scintillation, the twinkling of stars, which contributes a noise
proportional to the signal itself and therefore sets a hard ceiling per
exposure. No aperture photometry from the ground beats it in a single
frame; the way past it is stacking many frames, which is what the stack
planner quantifies.

\subsection{Extended sources: galaxies, nebulae, planets}

Set \textsf{Source geometry} to \code{extended} and fill in the source
area in arcsec$^2$. The convention to remember is that \emph{the
magnitude you enter remains the total, integrated magnitude of the
object} --- the number catalogues quote --- and the program spreads it
uniformly over the area you declare. Some example areas: a galaxy listed
as $2' \times 1'$ covers $\pi ab/4 \simeq 5\,650$ arcsec$^2$ if those are
full axes; Jupiter at opposition ($\simeq44''$ diameter) covers
$\simeq1\,500$ arcsec$^2$; the Ring Nebula ($\simeq 86'' \times 62''$)
about $4\,200$.

Physically the program then does three things differently from a point
source. First, your aperture no longer captures ``all the light minus
seeing losses'' but simply the fraction of the object's area it covers
--- an aperture smaller than the galaxy sees proportionally less of it.
Second, there is no seeing concentration: the brightest pixel holds the
mean surface brightness, not a stellar peak, which typically moves the
saturation limit by orders of magnitude (this is why you can expose far
longer on a galaxy than on a star of the same total magnitude). Third,
the S/N refers to the light inside your aperture, which is the quantity
your photometry software will actually measure. The approximation
declared in the model --- uniform brightness --- is good when the object
is much larger than the seeing disc; for objects comparable to the
seeing (small planetary nebulae), treat the result as indicative.

\subsection{Target S/N mode and the stack planner}

Typing a number into \textsf{Target S/N} reverses the question: instead
of ``what S/N in this exposure'' the program answers ``what exposure for
this S/N''. The inversion is done in closed form with the complete noise
budget --- including scintillation, which matters enormously here: for a
bright star a solver without scintillation would confidently prescribe a
short exposure that can never reach your target, no matter what the
photon statistics say. If the S/N you request lies above the
scintillation ceiling for a single frame, the required exposure grows
until the ceiling (which rises as $\sqrt{t}$) meets it.

The required exposure, however, frequently collides with saturation: a
bright target may need 300\,s for your S/N while its peak pixel saturates
in 20. This is exactly the situation the \textbf{stack planner} resolves.
After every run with a target S/N, the label under CALCULATION reports a
complete plan of the form:

\begin{quote}
\emph{Stack plan (aperture): 16 $\times$ 19 s (sub limited by saturation)
= 304 s total for S/N 500; read-noise penalty +2.1\% vs one ideal
exposure. Background beats read noise above 45 s/frame.}
\end{quote}

Reading it: the longest safe single frame is 19\,s (set by the brightest
pixel; or by your own preference if you typed one into \textsf{Stack
sub-exposure}); sixteen such frames reach the target S/N; splitting the
exposure costs only 2.1\% extra total time compared with one
hypothetical long frame, because each frame adds its own read noise; and
frames longer than 45\,s would be sky-limited, meaning read noise no
longer matters much beyond that length. If the penalty is large, your
frames are too short for your read noise --- raise the sub-exposure or
use a lower-noise gain setting.

\subsection{Differential photometry (variable stars, exoplanets)}

Photometric monitoring is almost never absolute: you measure your target
\emph{against} a comparison star in the same field, and the number that
matters is the scatter of that difference, in millimagnitudes. Enter the
comparison star's magnitude in \textsf{Comparison star mag} and the
program runs the identical calculation for it, combines the two error
budgets, and reports, in the PHOTOMETRY label and the outputs tab:

\begin{quote}
\emph{Differential precision vs m=11.5 comparison: 3.2 mmag per 60 s
frame (target S/N 520, comparison S/N 610).}
\end{quote}

Both stars' budgets include their own scintillation, treated as
uncorrelated between the two --- a conservative assumption at the
arcminute separations of typical comparison stars. To judge feasibility:
a typical exoplanet transit is 5--30 mmag deep, a bright-star transit
like HD 189733 about 25 mmag; with 3.2 mmag per frame you resolve such a
transit comfortably, and binning $N$ frames improves the precision by
$\sqrt{N}$. The AAVSO reporting convention likewise wants the per-frame
uncertainty --- this number is exactly that.
